\section{Evaluation}
\label{sec:eval}

We evaluate the return on verification investment across four strategies
operating at three budget levels (the fraction of outputs the judge may verify and intervene on).

\fakepara{Experimental Setup.}
We construct a balanced set of 200 queries: 100 knowable (verifiable facts
spanning multiple subjects) and 100 unknowable
(fabrication prompts for nonexistent people, fictional syndromes, future events,
and nonexistent citations). Ground truth labels are established by construction.
We run all queries on four model families: OLMo-3
7B~\cite{olmo2025olmo3}, Llama-3.1 8B~\cite{llama31}, Qwen3
4B~\cite{qwen3}, and Mistral 7B~\cite{mistral7b}, all instruction-tuned
for question-answering behavior.


\fakepara{Conditions.}
Our evaluation instantiates the runtime policy in a single-step setting:
the judge prioritizes outputs under a fixed verification budget,
verifies the selected outputs, and confirmed fabrications are replaced by abstention.
We compare four budget allocation strategies:
\begin{enumerate*}[label=(\roman*)]
\item \textbf{No judge.} Raw model output; the unverified baseline.
\item \textbf{Text-guided judge (length).} Verifies using response
word count, a text-channel signal needing no tensor interface.
\item \textbf{Tensor-guided judge.} Uses per-token entropy.
\item \textbf{Composed judge.} Tensor-guided selection for general queries; a
  substitute oracle for citation-like queries (see below).
\end{enumerate*}

The first three conditions differ in how outputs are prioritized for verification,
while the composed judge additionally demonstrates \emph{tool verification} for a known complementary failure mode.
We evaluate each condition at 10\%, 20\%, and 30\% verification budgets
and measure end-to-end accuracy.
Verification outcomes are simulated (1{,}000 seeded Monte Carlo trials):
a selected fabrication is caught with assumed probability 0.938 and a selected
knowable error is corrected with assumed probability 0.95. The former is the
stratified evaluator's agreement with a blinded author review of 80 items
(75/80), used as a simulation parameter rather than a measured catch rate.
Response correctness was labeled programmatically where the query admits it
(370 of 800) and otherwise by Qwen3-4B (430), which is also a model under
test; the labeling task differs from the generation task, and the 80-item
review above is the check on it.

\begin{figure}[t]%
\centering%
\includegraphics[width=\columnwidth]{figures/exp27_aggregate_budget_curve.pdf}%
\caption{End-to-end accuracy vs.\ verification budget for four judge conditions, with the 800 responses from four architectures pooled under one budget.}%
\label{fig:budget_curve}%
\end{figure}

\autoref{fig:budget_curve} shows that
\textbf{tensor-guided verification outperforms the text baseline
at every budget level}:
+3.3pp at 10\%, +4.6pp at 20\%, +3.3pp at 30\%.
We do not claim entropy is the better detector: length alone ranks at AUC
0.85--0.97 per family and exceeds entropy on two of four. The budget advantage
comes from which items each signal selects, and length, unlike entropy, is a
surface regularity training can move without changing the computation beneath
it. Concurrent work scales verification on the judge side by
taking the expectation over a judge's scoring-token distribution rather than its
discrete score~\cite{kwok2026llmasaverifiergeneralpurposeverificationframework};
our signals instead come from the generator's own computation, which---unlike a
judge's scores---the model does not separately control (\autoref{sec:threat}).
Per-model results, each with its own budget, show the same pattern.
At 10\% budget, tensor-guided verification improves accuracy over the
unverified baseline for every model tested: OLMo
(69.5\% $\to$ 75.1\%), Llama (82.5\% $\to$ 88.1\%), Qwen (87.0\%
$\to$ 91.7\%), and Mistral (64.0\% $\to$ 72.0\%).
The margin over the length baseline is not uniform, however: averaged
per family it is +2.6/+3.9/+2.6pp, and on Mistral at 30\% budget length
ranking wins by 2.2pp.
The signal is not confined to small local models: using only the top-5
log-probabilities that serverless APIs return (a lower bound on full-vocabulary
entropy), mean entropy separates knowable from unknowable queries on five further API
endpoints from 4B to 235B parameters (four further architectures, two of them
mixture-of-experts, plus a Mistral-7B repeat) with AUC
0.65--0.88~\cite{mason2026epistemic}.

\fakepara{Composed judges and complementary failure modes.}
At 30\% budget, the composed judge (91.5\%) outperforms the tensor-guided judge
alone (90.9\%). This is an upper bound rather than an achievable result: the
composed condition routes citation-like queries to a substitute oracle reading
the ground-truth label, so it measures the headroom a perfect lookup would buy,
not what a deployed one recovers.

\fakepara{Effect of post-training.}
To test whether standard post-training degrades the telemetry signals used by our interface, we run the same probes across four public OLMo-3 checkpoints---base, SFT, DPO, and RLVR---using greedy decoding on 16 knowable and 32 fabrication prompts.
Post-training reduces entropy for both classes, but does so disproportionately for knowable responses.
Mean entropy on knowable queries falls from 1.18 nats in the base model to 0.18 after post-training,
while entropy on the fabrications the model still produces falls from 1.49 to
0.98 after SFT and remains nearly unchanged through RLVR (0.99); taken over all
fabrication prompts, including the refusals, the RLVR figure is 0.92.
Consequently, the separation between knowable and fabricated responses \emph{increases},
with entropy-based discrimination improving from AUC 0.67 to 1.00 on these
probes. The ceiling depends on which knowable probes are counted: adding the
run's true-but-implausible category (\eg\ \autoref{fig:entropy_trace}) gives
0.63 rising to 0.81--0.87 instead, so the widening is robust while the perfect
ceiling is not. Private-context probes are excluded throughout, since trained
refusals collapse their entropy.
Post-training also changes the model's testimony: the fraction of fabrication prompts that elicit a refusal rises from 0\% to 56\%.
Thus, in this diagnostic, standard post-training improves the model's observable behavior while preserving---and in fact strengthening---the telemetry
signal that distinguishes the remaining fabrications from knowable responses.
